\section{模的正向极限}

记号约定:$\left(\left(A_{\alpha}\right)_{\alpha\in I},\left(\phi_{\beta\alpha}\right)_{\alpha,\beta\in I}\right)$是环的正系统,$A$是环.

\begin{definition}{左模的正系统和正向极限}{}
	设$\left(\left(E_{\alpha}\right)_{\alpha\in I},\left(f_{\beta\alpha}\right)_{\alpha,\beta\in I}\right)$是交换群的正系统.如果
	\begin{enumerate}
		\item 对每个$\alpha\in I$,$E_{\alpha}$是左$A_{\alpha}$-模，
		\item $\forall\alpha,\beta\in I\forall\lambda_{\alpha}\in A_{\alpha}\forall x_{\alpha}\in E_{\alpha}\left(\alpha\leq\beta\implies f_{\beta\alpha}\left(\lambda_{\alpha}x_{\alpha}\right)=\phi_{\beta\alpha}\left(\lambda_{\alpha}\right)f_{\beta\alpha}\left(x_{\alpha}\right)\right)$,
	\end{enumerate}
	则称$\left(\left(E_{\alpha}\right)_{\alpha\in I},\left(f_{\beta\alpha}\right)_{\alpha,\beta\in I}\right)$是左$A_{\alpha}$-模的逆系统,左$\varinjlim A_{\alpha}$-模$\varinjlim E_{\alpha}$是其正向极限.
\end{definition}

\begin{proposition}{模的正系统之间的同态的逆向极限}{}
	设$\left(\left(E_{\alpha}\right)_{\alpha\in I},\left(f_{\beta\alpha}\right)_{\alpha,\beta\in I}\right),\left(\left(F_{\alpha}\right)_{\alpha\in I},\left(g_{\beta\alpha}\right)_{\alpha,\beta\in I}\right)$是$A_{\alpha}$-模的正系统,$A=\varinjlim A_{\alpha}$,对每个$\alpha\in I$有
	\begin{align*}
		u_{\alpha}\in\Hom_{A_{\alpha}-\mathsf{Mod}}\left(E_{\alpha},F_{\alpha}\right),
	\end{align*}
	并且$\left(u_{\alpha}\right)_{\alpha\in I}$是$A_{\alpha}$-线性映射的正系统,则$\varinjlim u_{\alpha}\in\Hom_{A-\mathsf{Mod}}\left(\varinjlim E_{\alpha},\varinjlim F_{\alpha}\right)$.
\end{proposition}

\begin{proposition}{模的逆向极限的正合性}{}
	设$\left(\left(E_{\alpha}\right)_{\alpha\in I},\left(f_{\beta\alpha}\right)_{\alpha,\beta\in I}\right),\left(\left(F_{\alpha}\right)_{\alpha\in I},\left(g_{\beta\alpha}\right)_{\alpha,\beta\in I}\right),\left(\left(G_{\alpha}\right)_{\alpha\in I},\left(h_{\beta\alpha}\right)_{\alpha,\beta\in I}\right)$是$A_{\alpha}$-模的正系统,其中
	\begin{itemize}
		\item $\forall\alpha\in I\left(u_{\alpha}\in\Hom_{A_{\alpha}-\mathsf{Mod}}\left(E_{\alpha},F_{\alpha}\right)\wedge v_{\alpha}\in\Hom_{A_{\alpha}-\mathsf{Mod}}\left(F_{\alpha},G_{\alpha}\right)\right)$,
		\item $\left(u_{\alpha}\right)_{\alpha\in I},\left(v_{\alpha}\right)_{\alpha\in I}$是$A_{\alpha}$-线性映射的正系统.
	\end{itemize}
	如果对每个$\alpha\in I$,序列$E_{\alpha}\xlongrightarrow{u_{\alpha}}F_{\alpha}\xlongrightarrow{v_{\alpha}}G_{\alpha}$正合,则$\varinjlim E_{\alpha}\xlongrightarrow{\varinjlim u_{\alpha}}\varinjlim F_{\alpha}\xlongrightarrow{\varinjlim v_{\alpha}}\varinjlim G_{\alpha}$也正合.
\end{proposition}

\begin{proposition}{生成集的正向极限}{}
	设$\left(\left(E_{\alpha}\right)_{\alpha\in I},\left(f_{\beta\alpha}\right)_{\alpha,\beta\in I}\right)$是$A_{\alpha}$-模的正系统.若对每个$\alpha\in I$,$S_{\alpha}$是$E_{\alpha}$的生成集,则$\bigcup\limits_{\alpha\in I}f_{\alpha}\left(S_{\alpha}\right)$是$\varinjlim E_{\alpha}$的生成集.
\end{proposition}

\begin{proposition}{直和的正向极限}{}
	设$\left(\left(E_{\alpha}\right)_{\alpha\in I},\left(f_{\beta\alpha}\right)_{\alpha,\beta\in I}\right)$是$A_{\alpha}$-模的正系统.如果
	\begin{enumerate}
		\item 对每个$\alpha\in I$,$E_{\alpha}$是它的一族子模$\left(M_{\alpha}^{\lambda}\right)_{\lambda\in\Lambda}$的直和,
		\item $\forall\alpha,\beta\in I\forall\lambda\in\Lambda\left(\alpha\leq\beta\implies f_{\beta\alpha}\left(M_{\alpha}^{\lambda}\right)\subseteq M_{\beta}^{\lambda}\right)$,
	\end{enumerate}
	则$\varprojlim E_{\alpha}=\bigoplus\limits_{\lambda\in\Lambda}\varprojlim M_{\alpha}^{\lambda}$.
\end{proposition}

\begin{corollary}{自由集和基集的正向极限}{}
	设$\left(\left(E_{\alpha}\right)_{\alpha\in I},\left(f_{\beta\alpha}\right)_{\alpha,\beta\in I}\right)$是$A_{\alpha}$-模的正系统,$\left(P_{\alpha}\right)_{\alpha\in I}$是其子集的正系统.
	\begin{enumerate}
		\item 若对每个$\alpha\in I$,$P_{\alpha}$是$E_{\alpha}$的自由子集,则$\varprojlim P_{\alpha}$是$\varprojlim E_{\alpha}$的自由子集.
		\item 若对每个$\alpha\in I$,$P_{\alpha}$是$E_{\alpha}$的基集,则$\varprojlim P_{\alpha}$是$\varprojlim E_{\alpha}$的基集.
	\end{enumerate}
\end{corollary}

本节接下来的内容中,默认$\forall\alpha\in I\left(A_{\alpha}=A\right)\wedge\forall\beta,\gamma\in I\left(\phi_{\beta\gamma}=\id_{A}\right)$.

\begin{proposition}{}{}
	对每个$A$-模的正系统$\left(\left(E_{\alpha}\right)_{\alpha\in I},\left(f_{\beta\alpha}\right)_{\alpha,\beta\in I}\right)$和左$A$-模$F$,典范映射
	\begin{align*}
		d_{F}:\Hom_{A-\mathsf{Mod}}\left(\varinjlim E_{\alpha},F\right)\to\varprojlim\Hom_{A-\mathsf{Mod}}\left(E_{\alpha},F\right),u\mapsto\left(u\circ f_{\alpha}\right)_{\alpha\in I}
	\end{align*}
	是$\Z$-模同构.
\end{proposition}

\begin{corollary}{}{}
	设$\left(\left(E_{\alpha}\right)_{\alpha\in I},\left(f_{\beta\alpha}\right)_{\alpha,\beta\in I}\right)$是$A$-模的正系统,$F,G$是$A$-模,$v\in\Hom_{A-\mathsf{Mod}}\left(F,G\right)$,对每个$\alpha\in I$有$\bar{v}_{\alpha}=\Hom_{A}\left(\id_{E_{\alpha}},F\right)$.
	\begin{enumerate}
		\item $\left(\bar{v}_{\alpha}\right)_{\alpha\in I}$是$\Z$-线性映射的逆系统.
		\item 如下图表交换:
		\begin{center}
			\begin{tikzcd}
				{\Hom_{A}\left(\varinjlim E_{\alpha},F\right)} && {\varprojlim\Hom_{A}\left(E_{\alpha},F\right)} \\
				\\
				{\Hom_{A}\left(\varinjlim E_{\alpha},G\right)} && {\varprojlim \Hom_{A}\left(E_{\alpha},G\right)}
				\arrow["{d_{F}}", from=1-1, to=1-3]
				\arrow["{\Hom_{A}(\id_{\varprojlim E_{\alpha}},v)}"', from=1-1, to=3-1]
				\arrow["{\varprojlim\bar{v}_{\alpha}}", from=1-3, to=3-3]
				\arrow["{d_{G}}"', from=3-1, to=3-3]
			\end{tikzcd}
		\end{center}
	\end{enumerate}
\end{corollary}

\begin{corollary}{}{}
	$\left(\left(E_{\alpha}^{\vee}\right),\left(f_{\beta\alpha}^{\top}\right)_{\alpha,\beta\in I}\right)$是右$A$-模的逆系统,$\varprojlim E_{\alpha}^{\vee}$典范同构于$\left(\varinjlim E_{\alpha}\right)^{\vee}$.
\end{corollary}

\begin{remark}{}{}
	设$E$是左$A$-模,$\left(M_{\alpha}\right)_{\alpha\in I}$是$M$的一族递增子模,$E=\bigcup\limits_{\alpha\in I}M_{\alpha}$.考虑$\left(M_{\alpha}\right)_{\alpha\in I}$构成的正系统,那么$\varprojlim M_{\alpha}$典范同构于$E$.特别地,每一个左$A$-模都是他的全部有限生成子模的正向极限.
\end{remark}

















































































